\section{A constructive lower certificate over the entire origin plane}
\label{supp:constructive-global}

A positive global infimum can be converted into an explicit numerical certificate when opposite winning regions contain connected cores with uniform clearance and signed-density margins. Let $K_+,K_-$ be nonempty compact connected sets. Assume that almost everywhere on their open radius-$\rho$ neighborhoods,
\[
 d\geq\kappa_+>0\quad\hbox{on }K_++B(0,\rho),\qquad
 d\leq-\kappa_-<0\quad\hbox{on }K_-+B(0,\rho).
\]
Suppose that each core meets the other's convex hull, both cores are contained in $\overline B(p,M)$, and
\[
 B(p,\delta)\subseteq\operatorname{conv}(K_+)
                         \cap\operatorname{conv}(K_-),\qquad\delta>0.
\]
Write $\kappa=\min(\kappa_+,\kappa_-)$. The opposite signs make the open neighborhoods disjoint, so all opposite-core point pairs have distance at least $2\rho$, and $M\geq\rho$.

\paragraph{Proposition (constructive global floor).}
For every $C>M$,
\begin{equation}
 \inf_c\Delta(c)\geq\kappa\min\left\{
 \frac{\sqrt3\rho^3}{2(C+M)},\;
 4\rho\delta\left(1-\frac MC\right)\right\}>0.
 \label{supp:eq:constructive-floor}
\end{equation}
The choice $C=2M$ gives the simpler bound $\kappa\min\{\sqrt3\rho^3/(6M),2\rho\delta\}$. The two expressions in the minimum are equal at
\begin{equation}
 C_* =\frac{\sqrt3\rho^2+
             \sqrt{3\rho^4+256\delta^2M^2}}{16\delta}>M.
 \label{supp:eq:optimized-split}
\end{equation}
The first decreases and the second increases with $C$, so this choice maximizes their minimum.

\emph{Proof.}
First consider $c\notin K_+\cup K_-$. Each core's angular image is a nonempty compact connected subset of the circle. If these images were disjoint, they would be closed arcs, at least one of length strictly less than $\pi$. The positive cone generated by that arc is convex and contains the corresponding core's translated convex hull. An opposite-core point in that hull would then have a direction in both images, a contradiction. Thus a ray from $c$ meets both cores. If $c$ belongs to one core, take a ray through any point of the other core; the first core contributes its disk centered at $c$.

Let the two core points on the resulting ray be $c+r_iu_0$, with $r_i\geq0$, and put $D=\max_i r_i$. Their separation gives $D\geq2\rho$. Set
\[
 \alpha=\arcsin\left(\frac{\sqrt3\rho}{2D}\right)\leq\frac\pi3.
\]
Every ray making angle $|\phi|\leq\alpha$ with $u_0$ acquires signed mass at least $\kappa_i\rho^2/2$ from each core disk. To verify this when $r_i\leq\rho$, both the origin and the point at radial distance $\rho$ lie in the disk, because $r_i\leq2\rho\cos\phi$. The disk's convexity includes the intervening radial segment, whose mass is $\kappa_i\int_0^\rho r\,dr$. When $r_i\geq\rho$, the disk--ray chord has half-length
\[
 \ell_i=\sqrt{\rho^2-r_i^2\sin^2\phi}\geq\rho/2
\]
and midpoint radius $r_i\cos\phi\geq\rho/2$. The whole chord lies on the positive ray, since $(r_i\cos\phi)^2-\ell_i^2=r_i^2-\rho^2\geq0$. Its mass is $2\kappa_i r_i\cos\phi\,\ell_i\geq\kappa_i\rho^2/2$. Polar Fubini handles the almost-everywhere density assumptions, and boundary endpoints do not affect the integrals. Hence the cancellation identity gives
\begin{equation}
 \Delta(c)\geq\kappa\rho^2\alpha
             \geq\frac{\sqrt3\kappa\rho^3}{2D}.
 \label{supp:eq:compact-core-floor}
\end{equation}
For $\|c-p\|\leq C$, we have $D\leq C+M$, proving the first bound in~\eqref{supp:eq:constructive-floor}.

For $t=\|c-p\|>M$, strict separation from each core and connectedness imply that the core and its convex hull have the same angular image. Every ray crossing $B(p,\delta)$ therefore meets both cores, at radial distances $r_i\geq t-M$. Each core disk contributes mass at least $2\kappa_i\rho r_i$. This is the full chord integral if $r_i\geq\rho$. If $r_i<\rho$, its positive-ray integral is $\kappa_i(r_i+\rho)^2/2\geq2\kappa_i\rho r_i$. Integrating across the disk's subtended angle yields
\[
 \Delta(c)\geq4\kappa\rho(t-M)\arcsin(\delta/t)
 \geq4\kappa\rho\delta(1-M/t).
\]
For $t>C$, this proves the second bound. Equating the two expressions gives $8\delta(C^2-M^2)=\sqrt3\rho^2C$ and~\eqref{supp:eq:optimized-split}.\hfill$\square$

\paragraph{Certificates for the evaluated configurations.}
For $h=1/2$, use the eroded cores from Section~\ref{supp:exterior}, with $\rho=\delta=1/8$ and $M=43/40+\sqrt5/4$. They satisfy mutual hull intrusion at $(1,1/4)$ and its reflection. For $h=1$, use the exclusive backbones, with $\rho=1/20$, $\delta=1/8$ and $M=11/10+1/\sqrt2$; the barycentric identity~\eqref{supp:eq:barycentric} establishes mutual intrusion. At equal priors, $\kappa=5/(4\pi)$ for area sampling and $25/(24\pi)$ for radius sampling. Substitution into~\eqref{supp:eq:constructive-floor} at $C_*$ gives the exact whole-plane certificate expressions. Directed decimal interval arithmetic, with exact rational bounds on square-root endpoints and on Machin's formula for $\pi$, gives the following downward-rounded certificates:
\begin{center}
\begin{tabular}{ccc}
\toprule
$h$ & Area lower bound & Radius lower bound \\
\midrule
$1/2$ & $0.0002050842462841$ & $0.0001709035385701$ \\
$1$   & $0.0000119104106520$ & $0.0000099253422100$ \\
\bottomrule
\end{tabular}
\end{center}
All entries are loss probabilities valid over the entire origin plane. The exact expressions and enclosing arithmetic intervals accompany the calculation.

\paragraph{Construction throughout the disjoint interval.}
For $0<h\leq\sqrt{3/5}$, define
\begin{align*}
 q&=\sqrt{1+h^2/4},&p&=(1/2,h/2),\\
 \rho&=\min\{h/4,b-q,q-a\},\\
 \delta&=\min\{h/2-\rho,b-\rho-\|p\|\},\\
 M&=b-\rho+\|p\|.
\end{align*}
These quantities satisfy $\rho,\delta>0$: $a<1\leq q\leq\sqrt{23/20}<b$, $\rho\leq h/4$, and $b-\rho\geq q>\|p\|$, where $q^2-\|p\|^2=3/4$. Moreover $\rho<1/5$, so the eroded radial interval is nonempty. Set
\[
 K_0=\{x:a+\rho\leq\|x\|\leq b-\rho,\ x_2\geq\rho\},
 \qquad K_1=(1,h)-K_0.
\]
These compact connected cores have radius-$\rho$ open neighborhoods inside their respective supports. Their convex hulls are the corresponding outer disk caps. The point $(1,h/2)$ belongs to $K_0$ and to the interior of $\operatorname{conv}(K_1)$; reflection proves the reverse inclusion. The definitions of $\delta$ and $M$ give the common hull disk and extent bound. Because the supports are disjoint, the signed-density margins are
\[
 \kappa=\frac{\min(\pi_0,\pi_1)}H
 \quad\hbox{or}\quad
 \kappa=\frac{\min(\pi_0,\pi_1)}{\pi(b-a)b}
\]
for area or radius sampling, respectively. Thus~\eqref{supp:eq:constructive-floor} supplies an explicit positive global loss certificate throughout the entire disjoint positive-displacement interval.
